\documentclass[tikz,border=8pt]{standalone}
\usepackage{tikz}
\usepackage{amsmath}
\usepackage{amssymb}
\usepackage{pgfplots}
\pgfplotsset{compat=1.18}
\usepackage{ctex}
\usetikzlibrary{calc,positioning,arrows.meta,trees}

% BST 结构示意
% 根=8，左子树：3(1, 6(4,7))，右子树：10(_, 14(13,_))
% 中序遍历：1,3,4,6,7,8,10,13,14

\begin{document}
\begin{tikzpicture}[
  every node/.style={},
  level distance=1.4cm,
  level 1/.style={sibling distance=4.8cm},
  level 2/.style={sibling distance=2.4cm},
  level 3/.style={sibling distance=1.3cm},
  edge from parent/.style={draw, ->, >=Stealth, thick}
]

%% ── 标题 ──────────────────────────────────────────────────
\node[font=\large\bfseries] at (0, 1.4) {BST 结构示意};
\node[font=\small,gray] at (0, 0.95) {二叉搜索树：左子树 $<$ 根 $<$ 右子树};

%% ── BST 树 ────────────────────────────────────────────────
\node[draw,circle,minimum size=0.78cm,fill=blue!10,font=\bfseries] (root) at (0,0) {8}
  child {
    node[draw,circle,minimum size=0.78cm,fill=blue!10,font=\bfseries] (n3) {3}
    child {
      node[draw,circle,minimum size=0.78cm,fill=blue!10,font=\bfseries] (n1) {1}
    }
    child {
      node[draw,circle,minimum size=0.78cm,fill=blue!10,font=\bfseries] (n6) {6}
      child {
        node[draw,circle,minimum size=0.78cm,fill=blue!10,font=\bfseries] (n4) {4}
      }
      child {
        node[draw,circle,minimum size=0.78cm,fill=blue!10,font=\bfseries] (n7) {7}
      }
    }
  }
  child {
    node[draw,circle,minimum size=0.78cm,fill=blue!10,font=\bfseries] (n10) {10}
    child[missing] {}
    child {
      node[draw,circle,minimum size=0.78cm,fill=blue!10,font=\bfseries] (n14) {14}
      child {
        node[draw,circle,minimum size=0.78cm,fill=blue!10,font=\bfseries] (n13) {13}
      }
      child[missing] {}
    }
  };

%% ── BST 性质标注 ──────────────────────────────────────────
% 左子树标注
\node[draw=gray!50,fill=white,thin,rounded corners=2pt,
      font=\scriptsize,inner sep=3pt,align=center]
  at (-5.5,-0.5) {左子树\\$\{1,3,4,6,7\}$\\全部 $<$ 8};
\draw[->,>=Stealth,gray!60,thin] (-4.7,-0.5) -- (-3.0,-0.3);

% 右子树标注
\node[draw=gray!50,fill=white,thin,rounded corners=2pt,
      font=\scriptsize,inner sep=3pt,align=center]
  at (5.5,-0.5) {右子树\\$\{10,13,14\}$\\全部 $>$ 8};
\draw[->,>=Stealth,gray!60,thin] (4.7,-0.5) -- (3.0,-0.3);

%% ── 子树内部关系标注 ──────────────────────────────────────
\node[draw=gray!50,fill=white,thin,rounded corners=2pt,
      font=\tiny,inner sep=2pt] at (-3.2,-3.5)
  {$\{1\} < 3 < \{4,7\}$};
\node[draw=gray!50,fill=white,thin,rounded corners=2pt,
      font=\tiny,inner sep=2pt] at (3.5,-3.5)
  {$\{13\} < 14$};

%% ── 中序遍历结果 ──────────────────────────────────────────
\node[draw=gray!50,fill=green!5,thin,rounded corners=3pt,
      font=\small,inner sep=5pt,align=center]
  at (0,-6.5)
  {\textbf{中序遍历}（In-order）= 升序输出\\
   $1 \to 3 \to 4 \to 6 \to 7 \to 8 \to 10 \to 13 \to 14$};

%% ── 说明文字 ──────────────────────────────────────────────
\node[draw=gray!50,fill=white,thin,rounded corners=3pt,
      font=\scriptsize,inner sep=4pt,align=left]
  at (0,-7.9)
  {BST 关键性质：对任意节点 $v$，\\
   左子树所有节点 $< v < $ 右子树所有节点\\
   $\Rightarrow$ 中序遍历恰好产生升序序列};

\end{tikzpicture}
\end{document}
